La implementación del procesador debe adaptarse a los recursos y señales de la
placa donde va a funcionar. Este capítulo presenta esa adaptación para la
Basys 3, explica cómo se generan el reloj y el reset del sistema, y muestra los
resultados obtenidos al implementar el diseño en la FPGA.

\section{Plataforma física}

La implementación de referencia utiliza una \tech{Basys 3} con el dispositivo
\tech{XC7A35T-1CPG236C}. La placa proporciona el reloj de entrada, el botón de
reset y la conexión USB-UART que permite intercambiar datos con la PC.

El módulo superior de plataforma, \code{basys3_top}, conecta esas señales con
el sistema. La tabla~\ref{tab:pines} identifica los pines utilizados y la
función de cada uno.

\begin{table}[H]
  \centering
  \small
  \begin{tabular}{llll}
    \toprule
    \textbf{Puerto} & \textbf{Dirección} & \textbf{Pin} & \textbf{Conexión} \\
    \midrule
    \code{clk_100mhz_i} & entrada & W5 & oscilador de 100 MHz \\
    \code{reset_btn_i} & entrada & U18 & botón central BTNC; 1 pulsado,
    0 libre \\
    \code{uart_rx_i} & entrada & B18 & PC transmite; reposo 1 \\
    \code{uart_tx_o} & salida & A18 & FPGA transmite; reposo 1 \\
    \bottomrule
  \end{tabular}
  \caption{Pines de la Basys 3.}
  \label{tab:pines}
\end{table}

Dentro de \code{basys3_top} se encuentran el generador de reloj, el
acondicionador de reset y \code{system_top}, que contiene el procesador y los
bloques de comunicación. La señal de recepción UART se sincroniza dentro de
\code{uart_core}. Las conexiones de pines y las restricciones temporales se
especifican en los archivos XDC del proyecto.

\section{Clock y reset}

\subsection{Generación del reloj}

El oscilador de la placa entrega un reloj de \SI{100}{\mega\hertz}. A partir
de él, el Clocking Wizard utiliza un MMCM y un buffer global BUFG para generar
el reloj funcional de \SI{50}{\mega\hertz} que comparten la CPU, la UART y el
resto del sistema.

El generador entrega la señal \code{locked} cuando su salida se estabiliza.
Mientras el reloj todavía no está estable, el reset funcional permanece
activo. El botón físico también reinicia el generador de reloj, mientras que la
pérdida de \code{locked} solicita el reset funcional sin reiniciar el MMCM.

\subsection{Fuentes y sincronización del reset}

El sistema puede recibir una solicitud de reset desde el botón físico o desde
un comando UART. Además, la plataforma solicita reset cuando el generador de
reloj todavía no está estable. Estas fuentes se combinan antes de distribuir
el reset funcional:

\begin{codebox}
async_reset_request = physical_reset_request
                      OR clock_reset_request
                      OR uart_reset_request
reset = reset_sync(clk, async_reset_request)
\end{codebox}

El bloque \code{reset_sync} activa el reset de forma asíncrona y lo libera de
forma sincronizada con el reloj funcional mediante dos etapas. Los registros
del sistema aplican sus efectos de reset en el flanco de reloj.

Cuando la solicitud llega por UART, primero se transmite completa la respuesta
\code{RESET/ACCEPTED} y, después del bit de parada del último byte, se registra
la solicitud de reset. La secuencia detallada se encuentra en
\doc{architecture/platform/reset_integration.md}{Integración del reset}.

\section{Implementación y análisis temporal}

La implementación transforma el diseño descrito en RTL en una configuración
que puede cargarse en la FPGA. Para ello, Vivado realiza la síntesis, asigna los
recursos físicos y determina las conexiones entre ellos mediante colocación y
ruteo. Con el diseño resultante se analiza si las señales cumplen las
restricciones temporales y se genera el bitstream para programar la placa.

Se utilizó Vivado ML Standard 2025.2, Clocking Wizard 6.0 y el perfil de
memoria \code{REFERENCE}, con un reloj funcional de \SI{50}{\mega\hertz},
equivalente a un período de \SI{20}{\nano\second}.

\subsection{Resultados temporales}

El análisis temporal comprueba que los datos lleguen a los registros dentro de
los intervalos requeridos. La comprobación de \emph{setup} evalúa que lleguen
a tiempo para el flanco de captura, mientras que la de \emph{hold} verifica que
permanezcan estables durante el intervalo necesario después de ese flanco.

La tabla~\ref{tab:resultados-timing} resume los márgenes obtenidos. WNS y WHS
son los peores márgenes de setup y hold, respectivamente; TNS y THS acumulan
los márgenes negativos de cada tipo. WPWS indica el peor margen para el ancho
de pulso del reloj. Un margen positivo indica que la restricción correspondiente
se cumple en el peor caso analizado.

\begin{table}[H]
  \centering
  \small
  \begin{tabular}{lr}
    \toprule
    \textbf{Métrica} & \textbf{Resultado} \\
    \midrule
    WNS & \SI[explicit-sign=+]{1.070}{\nano\second} \\
    TNS & \SI{0}{\nano\second} \\
    WHS & \SI[explicit-sign=+]{0.026}{\nano\second} \\
    THS & \SI{0}{\nano\second} \\
    WPWS & \SI[explicit-sign=+]{3.000}{\nano\second} \\
    Clock funcional & \SI{50}{\mega\hertz} \\
    Frecuencia estimada por setup & $\approx \SI{52.8}{\mega\hertz}$ \\
    \bottomrule
  \end{tabular}
  \caption{Resultados de timing de la implementación \code{REFERENCE}.}
  \label{tab:resultados-timing}
\end{table}

Los márgenes de setup y hold son positivos, y sus sumas de márgenes negativos
son cero. Por lo tanto, los caminos analizados cumplen esas restricciones para
el reloj de \SI{50}{\mega\hertz}.

El camino con menor margen de setup parte del registro \code{EX/MEM}, en el bit
3 del registro destino, y llega al bit 10 del registro de estado global. Su
demora de datos es de \SI{18.845}{\nano\second}, de los cuales
\SI{13.388}{\nano\second} corresponden a ruteo. La diferencia entre las
llegadas del reloj a ambos extremos, denominada \emph{skew}, es de
\SI{-0.034}{\nano\second}.

A partir del margen de setup puede calcularse una frecuencia estimada:
\[
  f_{\mathrm{estimada}} =
  \frac{1}{\SI{20}{\nano\second}-\SI{1.070}{\nano\second}}
  \approx \SI{52.8}{\mega\hertz}.
\]
Esta estimación supone que se mantienen el ruteo, el skew y la incertidumbre
del reloj. Sirve como referencia para setup; los análisis de la implementación
se realizaron a \SI{50}{\mega\hertz}.

El camino con menor margen de hold va de \code{ID/EX} a \code{EX/MEM}, para el
bit 29 de la instrucción. En ese camino, la demora de datos es de
\SI{0.362}{\nano\second} y el skew es de \SI{0.264}{\nano\second}.

\subsection{Recursos utilizados}

Además de la temporización, la implementación permite conocer qué recursos de
la FPGA ocupa el diseño. La tabla~\ref{tab:recursos} resume los principales
resultados para el perfil \code{REFERENCE}.

\begin{table}[H]
  \centering
  \small
  \begin{tabular}{lrr}
    \toprule
    \textbf{Recurso} & \textbf{Uso} & \textbf{Porcentaje} \\
    \midrule
    LUT de slice & 3466 & 16,66\,\% \\
    Registros & 1893 & 4,55\,\% \\
    LUT como memoria & 700 & -- \\
    MMCM & 1 & -- \\
    BUFG & 2 & -- \\
    IOB & 4 & -- \\
    BRAM & 0 & 0\,\% \\
    DSP & 0 & 0\,\% \\
    \bottomrule
  \end{tabular}
  \caption{Utilización de recursos de la implementación \code{REFERENCE}.}
  \label{tab:recursos}
\end{table}

De las 700 LUT utilizadas como memoria, 176 corresponden a IMEM, 512 a DMEM y
12 a las FIFO de UART. El banco de registros utiliza 992 registros. En esta
implementación, IMEM y DMEM conservan lectura combinacional y escritura
síncrona.

\subsection{Comprobaciones adicionales}

Vivado no reportó hallazgos en la comprobación de reglas de diseño (DRC) ni
errores de ruteo. La comprobación \code{check_timing} tampoco registró puntos
internos sin restricción temporal, registros sin reloj, lazos combinacionales
ni latches. También se analizaron las restricciones de los caminos de
recepción y transmisión UART: ambas tuvieron cobertura del 100\,\% para una
cota de \SI{20}{\nano\second}, con márgenes respectivos de
\SI[explicit-sign=+]{17.166}{\nano\second} y
\SI[explicit-sign=+]{13.654}{\nano\second}.

Las excepciones temporales de reset se limitaron a las entradas asíncronas de
las dos etapas del sincronizador y no quedaron excepciones ignoradas. La
revisión de cruces de reloj reconoció ese sincronizador y la lógica que combina
las fuentes de reset. Los 32 avisos de síntesis corresponden a las marcas de
validez de DMEM implementadas en LUTRAM.

La implementación generó el archivo \code{tp3-reference.bit}, utilizado para
programar la FPGA. El bitstream y los reportes se encuentran en el
\evidencefile{tp3-issue-17-evidencia.zip}{paquete de evidencia}.

\section{Verificación de plataforma}

La verificación de plataforma comprueba que el sistema pueda configurarse,
implementarse y funcionar con las señales de la placa. También revisa que el
reloj, el reset y las restricciones temporales estén conectados de acuerdo con
el diseño. La tabla~\ref{tab:casos-plataforma} presenta los grupos utilizados
para estas comprobaciones.

\begin{table}[H]
  \centering
  \small
  \begin{tabularx}{\textwidth}{L{2.9cm}Yr}
    \toprule
    \textbf{Grupo} & \textbf{Tema} & \textbf{Alcance} \\
    \midrule
    \code{TEST-PLAT-001} & Perfiles, capacidades y anchos & 15 casos \\
    \code{TEST-PLAT-002} & Clocking Wizard, locked y reset & 10 casos \\
    \code{TEST-PLAT-003} & Composición RTL y dependencias & 1 caso \\
    \code{TEST-PLAT-004} & Constraints, recursos y timing & 1 caso \\
    \code{TEST-PLAT-005} & Aceptación con un bitstream & 151 reejecuciones \\
    \code{TEST-PLAT-006} & Correspondencia documental & 1 caso \\
    \bottomrule
  \end{tabularx}
  \caption{Grupos de verificación de plataforma.}
  \label{tab:casos-plataforma}
\end{table}

El manifiesto contiene 28 casos de plataforma. Por su parte,
\code{TEST-PLAT-005} selecciona 151 casos de otras áreas para volver a
ejecutarlos en placa; esas reejecuciones se contabilizan por separado.

Los grupos \code{TEST-PLAT-001} y \code{TEST-PLAT-002} comprueban las
capacidades de memoria y los perfiles, además del arranque, las fuentes de
reset y la pérdida y recuperación de \code{locked}.

La revisión \code{TEST-PLAT-003} comprueba los 51 bloques funcionales de la CPU,
sus conexiones y la composición del sistema, incluidos el MMCM y las dos etapas
de \code{reset_sync}. Para
\code{TEST-PLAT-004} se utilizan el bitstream y los reportes de implementación
y temporización presentados en este capítulo. \code{TEST-PLAT-006} revisa la
correspondencia de la documentación con la entrega.

Las pruebas de \code{TEST-PLAT-005} se describen en el
capítulo~\ref{chap:aceptacion-fisica}, junto con los resultados obtenidos y su
alcance.

\section{Resultado de aceptación}

Los resultados de implementación muestran que el diseño cumple las
restricciones analizadas a \SI{50}{\mega\hertz} y que el bitstream puede
utilizarse para la validación en placa. El análisis conservó la relación entre
el reloj de entrada y el reloj funcional, y las excepciones de reset se
limitaron a los caminos indicados anteriormente.
